% Full density, affine-substitution, convolution, gamma and normal proofs.
\section{Convolution, gamma and normal laws: proofs}
\label{sec:annex-p15-convolution-gamma-normal}

The source presentation is Durrett's fixed author version
\cite[pp.~33--34, 49--52]{DurrettPTEAuthor2019}.
Compact-interval integrals of the continuously differentiable functions
used below agree with their ordinary calculus integrals by the
simple-function bounds in Proof~\ref{annex:p15.fubini-failures}.
Improper endpoint calculations below are increasing-domain limits of
nonnegative Lebesgue integrals. The only changes of measure used in
substitutions are the affine ones proved here.

\begin{annexitem}{annex:p15.density-integrals}{thm:p15-density-integrals}{Weighted measures and the meaning of a density}
\begin{proof}
The earlier nonnegative integral theorem proves that
\(A\mapsto\int_A f\,d\lambda^d\) is a measure. Its total
mass is one by assumption, so it is a probability law. For
\(h=\mathbf1_A\), the claimed integral identity is its definition.
Finite nonnegative combinations give the identity for simple \(h\).
For general nonnegative Borel \(h\), choose nonnegative simple
\(h_n\uparrow h\). Then \(h_nf\uparrow hf\) under the
nonnegative extended product convention: where \(f=0\) every
product is zero, where \(0<f<\infty\) it is ordinary scaling,
and where \(f=\infty\) the limit is zero if \(h=0\) and
infinity otherwise. Monotone convergence on both measure spaces
proves the identity. Applying it to \(|h|\) gives the integrability
criterion. Applying it separately to the two parts of an integrable
signed \(h\) and subtracting finite numbers gives the signed formula.

Every density has finite integral, so it is finite outside its
measurable infinite-value null set by the finite-integral criterion.
Setting it to zero there is a Borel replacement with the same
integrals. More generally measurable a.e. equality of nonnegative
functions preserves their integrals on every Borel set, hence their
density laws, by the established a.e. invariance theorem.

If \(f,g\) represent the same probability law, first make the
finite Borel replacements just described on the union of their two
null sets. They are now finite integrable real functions. For
\(E_n=\{f\geq g+1/n\}\), equality of the two laws gives
\(\int_{E_n}(f-g)\,d\lambda^d=0\). But the integrand on
\(E_n\) is at least \(1/n\), so
\(0\geq\lambda^d(E_n)/n\). Thus \(\lambda^d(E_n)=0\).
The union of these sets is \(\{f>g\}\), which is null.
Interchanging \(f,g\) proves the same for \(\{g>f\}\).
This gives a.e. uniqueness, including the original null-set versions.

A singleton in \(\mathbb R^d\) has Lebesgue measure zero by
the finite product rectangle formula and the one-dimensional
singleton result. Its integral against any density is zero.
The Dirac assignment is a probability measure because in a disjoint
union at most one member can contain \(c\); its value on a
singleton \(\{c\}\) is one, so it cannot have a Lebesgue
density. A variable equal to \(c\) a.s. has exactly those
event probabilities, proving the stated law interpretation.
\end{proof}
\end{annexitem}

\begin{annexitem}{annex:p15.affine-lebesgue}{thm:p15-affine-lebesgue}{Affine substitution proved from interval lengths}
\begin{proof}
The map \(T(x)=cx+d\), \(c\ne0\), and its inverse are
continuous, hence Borel measurable. In particular \(T(A)\)
is Borel whenever \(A\) is Borel, since it is the inverse
image of \(A\) under the continuous inverse map. The assignment
\(\eta(A)=\lambda(T(A))/|c|\) is a Borel measure: the
bijection preserves disjointness and unions. For real \(r<s\),
the image of \((r,s]\) is \((cr+d,cs+d]\) if \(c>0\),
and \([cs+d,cr+d)\) if \(c<0\). Either has length
\(|c|(s-r)\), because endpoint singletons are null.
Thus \(\eta((r,s])=s-r\). Uniqueness of the Borel length
measure in Theorem~\ref{thm:p15-lebesgue-unit-interval} gives
\(\eta=\lambda\), proving the measure-scaling formula.

For an indicator \(h=\mathbf1_A\), the integral of \(h(T(x))\)
is \(\lambda(T^{-1}(A))=\lambda(A)/|c|\), by the formula
applied to the inverse affine map. Finite simple combinations and
then increasing simple approximation prove the displayed integral
identity for every nonnegative Borel \(h\). The absolute-value
version proves equivalence of integrability, and the two signed
parts then give the integrable signed identity. Multiplying by
an interval indicator proves any restricted-domain version used below.

For the density transformation, the law of \(cX+d\) on a Borel
set \(A\) is
\(\int\mathbf1_A(cx+d)f(x)\,d\lambda(x)\).
Apply the proved substitution to
\(h(u)=\mathbf1_A(u)f((u-d)/c)\), a nonnegative Borel function.
The last integral becomes
\[
 \frac1{|c|}\int_A f((u-d)/c)\,d\lambda(u),
\]
which is the asserted density formula on every Borel set. If
\(c=0\), the transformed variable equals \(d\) everywhere,
and its law is \(\delta_d\); neither an inverse map nor a
division by zero is used.
\end{proof}
\end{annexitem}

\begin{annexitem}{annex:p15.convolution-law}{thm:p15-convolution-law}{Addition of independent coordinates and the cumulative formula}
\begin{proof}
Addition \(s(x,y)=x+y\) is continuous and Borel measurable.
The pushforward of the product probability under this map is
therefore a probability measure, proving existence of \(\mu*\nu\).
If independent \(X,Y\) have these laws, their joint law is
\(\mu\times\nu\) by Theorem~\ref{thm:p15-independent-product-law}.
Taking its pushforward under addition identifies the law of
\(X+Y\) with the defined convolution. The law-integration theorem
and Tonelli give the stated nonnegative integral formula.

For a fixed real \(z\), choose \(h=\mathbf1_{(-\infty,z]}\).
At fixed \(y\), the inner integral is
\[
 \int\mathbf1_{\{x+y\leq z\}}\,d\mu(x)
    =\mu((-\infty,z-y])=F(z-y).
\]
The other order similarly gives \(G(z-x)\). Tonelli proves
measurability of these section integrals and equality of the two
outer integrals. This proves the source's formula with \(dG\)
meaning the measure \(\nu\), not a derivative.

Commutativity follows by interchanging the two coordinates in the
nonnegative formula and using \(x+y=y+x\). For associativity,
let \(\eta\) be a third Borel probability. On a Borel set \(A\),
both \((\mu*\nu)*\eta\) and \(\mu*(\nu*\eta)\) have value
\[
 \int\!\int\!\int\mathbf1_A(x+y+z)
                   \,d\mu(x)\,d\nu(y)\,d\eta(z).
\]
To obtain the first expression, apply the two-factor formula to
the outer convolution and then to \(\mu*\nu\) at each fixed
\(z\). To obtain the second, first integrate over the other
grouping and then use Tonelli to put the three nonnegative integrals
in the displayed order. The finite product theorem identifies the
two grouped product domains; associativity of real addition gives
the same integrand. This proves equality on every Borel set.

Integration against \(\delta_c\) evaluates a nonnegative Borel
function at \(c\): this is first immediate for indicators, then
for simple functions and their increasing limits. Thus
\((\mu*\delta_c)(A)=\int\mathbf1_A(x+c)\,d\mu(x)=\mu(A-c)\).
Taking \(c=0\) proves the identity assertion.

For independent integer-valued \(X,Y\), the event \(X+Y=n\)
is the disjoint countable union of events \(\{X=k,Y=n-k\}\),
\(k\in\mathbb Z\). Independence factors each event probability.
Countable additivity gives the asserted nonnegative series, with
no interchange of a conditionally convergent signed sum.
\end{proof}
\end{annexitem}

\begin{annexitem}{annex:p15.convolution-density}{thm:p15-convolution-density}{The translated density and Tonelli's theorem}
\begin{proof}
The map \((z,y)\mapsto z-y\) is continuous, so
\((z,y)\mapsto f(z-y)\) is nonnegative product measurable.
The factor measures \(\lambda,\nu\) are sigma-finite. Tonelli
therefore makes \(h(z)=\int f(z-y)\,d\nu(y)\) a nonnegative
Borel measurable function. For any Borel \(A\), Tonelli and
the affine substitution at each fixed \(y\) give
\[
 \begin{aligned}
 \int_A h(z)\,d\lambda(z)
   &=\int\left(\int_A f(z-y)\,d\lambda(z)\right)d\nu(y)\\
   &=\int\left(\int\mathbf1_A(x+y)f(x)\,d\lambda(x)\right)d\nu(y)\\
   &=\int\left(\int\mathbf1_A(x+y)\,d\mu(x)\right)d\nu(y)
     =(\mu*\nu)(A).
 \end{aligned}
\]
The third equality is the proved density-integration formula, and
the last is the convolution law. Thus \(h\) is a density on
all Borel sets. Taking \(A=\mathbb R\) proves its integral
is one and hence that it is finite a.e.

For the source's cumulative-function route, specialize this same
calculation to \(A=(-\infty,z]\). The affine substitution
identifies \(F(z-y)\) with
\(\int_{(-\infty,z]}f(x-y)\,d\lambda(x)\); Tonelli then
puts the outer \(dG(y)\) integral inside and gives the density
formula exactly. This specialization needs no differentiability
of the cumulative function.

If \(\nu\) has density \(g\), apply the nonnegative density
formula to the Borel function \(y\mapsto f(z-y)\) for each
fixed \(z\). It gives
\(h(z)=\int f(z-y)g(y)\,d\lambda(y)\), including possible
infinite values. The nonnegative product convention handles any
zero/infinite density values. The already proved unit total
integral shows these exceptional infinite values form a null set.
\end{proof}
\end{annexitem}

\begin{annexitem}{annex:p15.factorized-joint-density}{thm:p15-factorized-joint-density}{Normalizing separated factors without assuming they were densities}
\begin{proof}
For any nonnegative measurable \(g\) and constant \(C\in[0,\infty]\),
\(\int Cg=C\int g\). The finite-constant cases are the
nonnegative scaling law. For \(C=\infty\), use monotone
convergence on \(r g\), \(r\uparrow\infty\); the limit
integral is zero if \(\int g=0\) and infinity otherwise,
as the convention requires. Repeated Tonelli, integrating one
coordinate at a time and using this constant identity, yields
\[
 \int_{\mathbb R^n}\prod_{i=1}^n g_i(x_i)\,d\lambda^n
                     =\prod_{i=1}^n\int_{\mathbb R}g_i\,d\lambda
                     =\prod_{i=1}^n c_i.
\]
The same calculation applies on any measurable rectangle. The
product integrand is measurable; for extended factors it is the
increasing limit of the products of their finite truncations,
with the prescribed zero-times-infinity convention.

Because the left side is the total mass of a density, it equals
one. No \(c_i\) can be zero, since then the finite product
would be zero; with all factors positive, none can be infinite,
since then the product would be infinite. Thus every \(c_i\)
is finite and positive, and \(g_i/c_i\) is a density.
For Borel \(A_i\), the rectangle calculation gives
\[
 \mathbb P(X_i\in A_i\text{ for all }i)
   =\prod_i\int_{A_i}g_i\,d\lambda
   =\prod_i\int_{A_i}(g_i/c_i)\,d\lambda,
\]
using \(\prod_i c_i=1\). Taking all sets but one to be the
whole line identifies the respective marginal densities. The
rectangle factorization is therefore precisely independence of
the variables. This includes \(n=1\), when the conclusion is
normalization of its one factor.
\end{proof}
\end{annexitem}

\begin{annexitem}{annex:p15.gamma-normalization}{thm:p15-gamma-normalization}{Gamma finiteness, recurrence and all stated moments}
\begin{proof}
For \(0<x<1\), \(x^{\alpha-1}e^{-x}\leq x^{\alpha-1}\),
whose integral is \(1/\alpha\) when \(\alpha>0\).
To justify this value within the present integral, integrate its
ordinary power antiderivative on \([\varepsilon,1]\), then
let \(\varepsilon\downarrow0\) by monotone convergence.
For \(x\geq1\), if \(\alpha\leq1\), the gamma integrand
is at most \(e^{-x}\). If \(\alpha>1\), put \(r=\alpha-1>0\).
The derivative of \(x^r e^{-x/2}\) has the sign of \(r-x/2\),
so this function is bounded on \([1,\infty)\) by a finite
constant \(M\). Consequently
\(x^{\alpha-1}e^{-x}\leq M e^{-x/2}\), an integrable bound.
This proves finiteness at both ends. Positivity follows because
the integrand is strictly positive on, for example, \((1,2)\),
which has positive measure, and a nonnegative function has zero
integral only when it vanishes a.e.

The same derivative argument proves, for any \(r>0,c>0\), that
\(x^r e^{-cx/2}\) is bounded for \(x\geq1\). Hence
\(x^r e^{-cx}\to0\). We will also use this elementary
exponential-tail limit in the normal-moment calculation.

Integration by parts on \([\varepsilon,R]\), where the
functions are continuously differentiable, gives
\[
 \int_\varepsilon^R x^\alpha e^{-x}\,d\lambda(x)
   =[-x^\alpha e^{-x}]_\varepsilon^R
       +\alpha\int_\varepsilon^R x^{\alpha-1}e^{-x}\,d\lambda(x).
\]
The boundary terms vanish as \(\varepsilon\downarrow0\),
\(R\uparrow\infty\), using \(\alpha>0\) and the tail
limit. Monotone convergence on both nonnegative integrals gives
\(\Gamma(\alpha+1)=\alpha\Gamma(\alpha)\).
The exponential antiderivative gives \(\Gamma(1)=1\), and
induction yields \(\Gamma(k+1)=k!\).

Use the proved positive affine scaling with \(u=\lambda x\).
For every \(s>0\),
\[
 \int_0^\infty x^{s-1}e^{-\lambda x}\,d\lambda(x)
                         =\lambda^{-s}\Gamma(s).
\]
Here the scalar \(\lambda^{-s}\) comes from
\(x^{s-1}=(u/\lambda)^{s-1}\) and the measure scaling
factor \(1/\lambda\). Choosing \(s=\alpha\) proves
that the displayed gamma density integrates to one. Choosing
\(s=\alpha+k\) and using the density-integration theorem gives
\[
 \mathbb E[X^k]
  =\frac{\lambda^\alpha}{\Gamma(\alpha)}
       \lambda^{-(\alpha+k)}\Gamma(\alpha+k)
  =\frac{\Gamma(\alpha+k)}{\lambda^k\Gamma(\alpha)}.
\]
The zeroth moment is the integral of the constant one; the
positive integer moments are finite by the finiteness just proved.
For \(k=1,2\), the recurrence yields mean \(\alpha/\lambda\)
and second moment \(\alpha(\alpha+1)/\lambda^2\).
Subtracting the squared mean gives variance \(\alpha/\lambda^2\).
Putting \(\alpha=1\) gives the stated exponential density,
factorial moments and their mean/variance specializations.
\end{proof}
\end{annexitem}

\begin{annexitem}{annex:p15.beta-gamma-convolution}{thm:p15-beta-gamma-convolution}{Both source beta-gamma arguments and the gamma sum law}
\begin{proof}
First establish the beta integral without assuming its gamma
identity. On \((0,1/2]\), the factor \((1-u)^{\beta-1}\)
is bounded by \(\max(1,2^{1-\beta})\), while
\(u^{\alpha-1}\) is integrable at zero for \(\alpha>0\).
On \([1/2,1)\), the factor \(u^{\alpha-1}\) is bounded
by \(\max(1,2^{1-\alpha})\), and the other factor is
integrable at one by the proved reflection substitution. Thus
\(B(\alpha,\beta)<\infty\). It is positive because its
integrand is positive on the whole open interval. Substitution
\(v=1-u\), with absolute scaling factor one, gives
\(B(\alpha,\beta)=B(\beta,\alpha)\).
The change \(y=zu\), \(z>0\), similarly gives the exact
source rescaling identity
\[
 \begin{aligned}
 \int_0^z(z-y)^{\alpha-1}y^{\beta-1}\,d\lambda(y)
  &=z^{\alpha+\beta-1}
               \int_0^1(1-u)^{\alpha-1}u^{\beta-1}\,d\lambda(u)\\
  &=z^{\alpha+\beta-1}B(\beta,\alpha).
 \end{aligned}
\]
Each is a one-variable affine substitution; all integrands are
nonnegative and endpoint values can be set to zero.

\paragraph{The normalization proof.}
The two gamma densities have already been normalized independently
of any beta-gamma identity. Independent variables with those laws
exist on their finite product probability space. The convolution
density theorem therefore gives a density of total integral one.
For \(z>0\), the overlap of their positive supports is \(0<y<z\),
and direct substitution of the densities gives
\[
 \begin{aligned}
 h(z)
  &=\frac{\lambda^{\alpha+\beta}e^{-\lambda z}}
           {\Gamma(\alpha)\Gamma(\beta)}
            \int_0^z(z-y)^{\alpha-1}y^{\beta-1}\,d\lambda(y)\\
  &=c_{\alpha,\beta}\lambda^{\alpha+\beta}
                     z^{\alpha+\beta-1}e^{-\lambda z},
 \qquad
 c_{\alpha,\beta}=\frac{B(\beta,\alpha)}
                         {\Gamma(\alpha)\Gamma(\beta)}.
 \end{aligned}
\]
For \(z\leq0\) the density is zero. The gamma integral
calculation at shape \(\alpha+\beta\) gives
\(\int_0^\infty\lambda^{\alpha+\beta}z^{\alpha+\beta-1}
e^{-\lambda z}\,d\lambda(z)=\Gamma(\alpha+\beta)\).
Since \(h\) has integral one, its finite positive constant must
satisfy \(c_{\alpha,\beta}\Gamma(\alpha+\beta)=1\).
It follows that
\[
 B(\beta,\alpha)=
       \frac{\Gamma(\alpha)\Gamma(\beta)}{\Gamma(\alpha+\beta)}
\]
and that \(h=g_{\alpha+\beta,\lambda}\). Symmetry of the
beta integral gives the displayed identity in the body's parameter
order. This proof normalizes the convolution using a gamma integral
already known to be finite; it has not assumed the identity being proved.

\paragraph{The nonnegative double-integral proof.}
The source gives a second route to the same identity. Multiply the
rescaling identity by \(e^{-z}\) and integrate over \(z>0\).
The left integral after rescaling is
\(\Gamma(\alpha+\beta)B(\beta,\alpha)\).
The other expression is an integral over the strict triangle
\(0<y<z\) of a nonnegative Borel function. Tonelli changes its
order and the affine translation \(t=z-y\) gives
\[
 \begin{aligned}
 &\int_0^\infty\int_0^z
      (z-y)^{\alpha-1}y^{\beta-1}e^{-z}\,d\lambda(y)\,d\lambda(z)\\
 &\quad=\int_0^\infty y^{\beta-1}e^{-y}
       \left(\int_y^\infty(z-y)^{\alpha-1}e^{-(z-y)}
                                      \,d\lambda(z)\right)d\lambda(y)\\
 &\quad=\Gamma(\alpha)\int_0^\infty y^{\beta-1}e^{-y}\,d\lambda(y)
       =\Gamma(\alpha)\Gamma(\beta).
 \end{aligned}
\]
Nonnegativity justifies the interchange before finiteness is known;
the final finite values establish it. Division by the positive
\(\Gamma(\alpha+\beta)\) proves the same beta-gamma identity
without using the normalization argument. This completes both
source routes.

For a finite mutually independent family with common rate,
the sum of any initial block is a measurable function of that
block and is independent of the next variable by the disjoint-group
theorem. Induction on the family size therefore adds all shapes.
The case of \(n\) exponentials uses shape one for every variable,
so the resulting shape is \(n\).

The beta density integrates to one directly from its finite positive
normalizer. Its first and second moments are
\(B(\alpha+1,\beta)/B(\alpha,\beta)\) and
\(B(\alpha+2,\beta)/B(\alpha,\beta)\). Using the proved
identity and the gamma recurrence gives
\[
 \mathbb E U=\frac\alpha{\alpha+\beta},\qquad
 \mathbb E U^2=\frac{\alpha(\alpha+1)}
                  {(\alpha+\beta)(\alpha+\beta+1)}.
\]
Subtracting the squared first moment gives the asserted variance.
All these moments are finite also because the beta variable lies
in \((0,1)\) a.s.
\end{proof}
\end{annexitem}

\begin{annexitem}{annex:p15.normal-normalization}{thm:p15-normal-normalization}{Gaussian normalization from affine substitutions and Tonelli}
\begin{proof}
Put \(J=\int_0^\infty e^{-x^2}\,d\lambda(x)\).
It is positive, and it is finite: the integrand is at most one
on \((0,1)\), and at most \(e^{-x}\) when \(x\geq1\).
Tonelli gives
\(J^2=\int_0^\infty\int_0^\infty e^{-x^2-y^2}
\,d\lambda(y)\,d\lambda(x)\).
For each fixed \(x>0\), apply only the one-variable affine
scaling \(y=xt\). Then apply Tonelli again:
\[
 \begin{aligned}
 J^2
  &=\int_0^\infty\int_0^\infty
       x e^{-(1+t^2)x^2}\,d\lambda(t)\,d\lambda(x)\\
  &=\int_0^\infty\left(\int_0^\infty
       x e^{-(1+t^2)x^2}\,d\lambda(x)\right)d\lambda(t).
 \end{aligned}
\]
For fixed \(c>0\), the antiderivative
\(-e^{-cx^2}/(2c)\) on bounded positive intervals, followed
by monotone convergence to \((0,\infty)\), gives
\(\int_0^\infty x e^{-cx^2}\,d\lambda(x)=1/(2c)\).
Taking \(c=1+t^2\) yields
\[
 J^2=\frac12\int_0^\infty\frac{d\lambda(t)}{1+t^2}
     =\frac12[\arctan t]_0^\infty=\frac\pi4.
\]
Thus \(J=\sqrt\pi/2\). Reflection doubles this half-line
integral, and affine scaling by \(\sqrt2\) gives
\(\int_{\mathbb R}e^{-x^2/2}\,d\lambda(x)=\sqrt{2\pi}\).
This argument used one-variable substitutions at fixed parameters
and two applications of Tonelli, not a two-variable change-of-variables
formula.

For the standard normal density, the same antiderivative at
\(c=1/2\) gives
\[
 \mathbb E|Z|=\frac2{\sqrt{2\pi}}
                  \int_0^\infty x e^{-x^2/2}\,d\lambda(x)
             =\sqrt{2/\pi}<\infty.
\]
Reflection makes the integral of the odd function
\(x e^{-x^2/2}\) equal its negative. It is absolutely integrable
by the displayed calculation, so its integral is zero and
\(\mathbb E Z=0\).
For the second moment, integration by parts on \([-R,R]\)
gives
\[
 \int_{-R}^R x^2 e^{-x^2/2}\,d\lambda(x)
   =\int_{-R}^R e^{-x^2/2}\,d\lambda(x)-2R e^{-R^2/2}.
\]
The last term tends to zero: for \(R\geq1\) it is bounded
by \(2R e^{-R/2}\), which tends to zero by the exponential-tail
bound already proved. Monotone convergence on the two nonnegative
integrals therefore gives a second-moment integral equal to
\(\sqrt{2\pi}\). Dividing by the normalizer proves
\(\mathbb E Z^2=1\) and hence variance one.

For \(a>0\), the affine image \(\mu+\sqrt a Z\) has the
displayed density \(\phi_{\mu,a}\), by the proved density
transformation. Finite absolute first and second moments follow
from those of \(Z\), for example from
\((\mu+\sqrt a Z)^2\leq2\mu^2+2aZ^2\).
The previously proved affine expectation and variance identities
give mean \(\mu\) and variance \(a\). Thus the density
normalization and the interpretation of both parameters are established.
For the declared \(a=0\) law, the variable is constant \(\mu\)
a.s., which directly gives the same mean and variance assertions.
\end{proof}
\end{annexitem}

\begin{annexitem}{annex:p15.normal-sums}{thm:p15-normal-sums}{The complete square calculation and all variance cases}
\begin{proof}
For \(a>0\), \(c\ne0\), the affine density formula for
\(cX+d\) is
\[
 \frac1{|c|\sqrt{2\pi a}}
       \exp\left(-\frac{((t-d)/c-\mu)^2}{2a}\right)
  =\frac1{\sqrt{2\pi c^2a}}
       \exp\left(-\frac{(t-(c\mu+d))^2}{2c^2a}\right).
\]
This is \(\operatorname{normal}(c\mu+d,c^2a)\), including
negative \(c\). If \(c=0\) the variable is constant \(d\),
and if \(a=0\) it equals \(c\mu+d\) a.s.; these cases give
the declared zero-variance normal laws separately.

First prove the independent sum formula for zero means and positive
variances \(a,b\). By the density convolution theorem, its
density at a real \(z\) is
\[
 \frac1{2\pi\sqrt{ab}}
       \int_{\mathbb R}
          \exp\left(-\frac{x^2}{2a}-\frac{(z-x)^2}{2b}\right)
                                             \,d\lambda(x).
\]
Complete the square by the explicit identity
\[
 \begin{aligned}
 bx^2+a(z-x)^2
    &=(a+b)x^2-2azx+az^2\\
    &=(a+b)\left(x-\frac{az}{a+b}\right)^2
                          +\frac{ab}{a+b}z^2.
 \end{aligned}
\]
After division by \(2ab\), the exponential splits into the
constant factor \(e^{-z^2/(2(a+b))}\) and the centered Gaussian
kernel with center \(az/(a+b)\) and variance
\(v=ab/(a+b)>0\). The proved Gaussian normalization and affine
translation give
\[
 \int_{\mathbb R}
       \exp\left(-\frac{(x-az/(a+b))^2}{2v}\right)d\lambda(x)
                         =\sqrt{2\pi v}
                         =\sqrt{\frac{2\pi ab}{a+b}}.
\]
Restoring the original prefactor, the convolution density is
\[
 \frac1{2\pi\sqrt{ab}}\sqrt{\frac{2\pi ab}{a+b}}
              e^{-z^2/(2(a+b))}
       =\frac1{\sqrt{2\pi(a+b)}}e^{-z^2/(2(a+b))}.
\]
This is the zero-mean normal density of variance \(a+b\),
with every constant accounted for.

For general means, \(X-\mu\) and \(Y-\nu\) are independent
by measurable transformation and have the centered laws just
considered. Their sum has mean parameter zero and variance parameter
\(a+b\). Shifting back by \(\mu+\nu\) gives the required
law for \(X+Y\).

If \(a=0\), then \(X=\mu\) a.s., so \(X+Y=\mu+Y\)
a.s. The latter has law \(\operatorname{normal}(\mu+\nu,b)\)
by the affine result. Measurable a.s. equality preserves the law,
so this proves the case, including \(b=0\). The case \(b=0\)
is the same argument with the variables interchanged. No positive-
variance density is evaluated at a zero variance.
Finally, for a finite mutually independent family, an initial
partial sum is independent of the next variable by the disjoint-
group theorem. Induction applies the proved two-variable formula
and adds all mean and variance parameters, including zeros.
\end{proof}
\end{annexitem}

\begin{annexitem}{annex:p15.convolution-calculations-boundaries}{thm:p15-convolution-calculations-boundaries}{Worked densities and failed independence or parameters}
\begin{proof}
For the elementary binary example, the integer convolution formula
gives masses \((1/2)^2=1/4\) at zero and two and the sum of
two such products, \(1/2\), at one. Other integers have zero
mass. For independent uniform variables, the density convolution
integrand is the indicator of \((0,1)\cap(z-1,z)\).
This intersection is empty when \(z\leq0\) or \(z\geq2\),
has length \(z\) when \(0<z\leq1\), and length \(2-z\)
when \(1<z<2\). Its total integral is
\(\int_0^1z\,d\lambda(z)+\int_1^2(2-z)\,d\lambda(z)=1\).
Every point mass is zero by the density theorem. These verify
all statements of Example~\ref{def:p15-convolution-elementary-examples}.

At rate two, the sum of two independent exponential variables has
gamma shape two by the gamma-sum theorem. Since \(\Gamma(2)=1\),
its density is \(4z e^{-2z}\), its mean is \(2/2=1\), and
its variance is \(2/2^2=1/2\). For beta parameters \(2,3\),
the beta-gamma identity gives
\(B(2,3)=\Gamma(2)\Gamma(3)/\Gamma(5)=1\cdot2/24=1/12\),
so the density is \(12u(1-u)^2\) on \((0,1)\).

For the normal example, the sum of the mean parameters is
\(1+(-2)=-1\), and the sum of the variance parameters is
\(4+9=13\). Substitution in the normal density formula gives
exactly \((26\pi)^{-1/2}e^{-(z+1)^2/26}\).

For rates one and two, the independent exponential convolution at
\(z>0\) gives
\[
 \int_0^z e^{-(z-y)}2e^{-2y}\,d\lambda(y)
       =2e^{-z}\int_0^z e^{-y}\,d\lambda(y)
       =2(e^{-z}-e^{-2z}).
\]
It is zero for \(z\leq0\). Independent means and variances
add, so the mean is \(1+1/2=3/2\) and the variance is
\(1+1/4=5/4\). A gamma law of shape two with this mean
would have rate \(\lambda=4/3\). Its variance would then
be \(2/(4/3)^2=9/8\), which differs from \(5/4\).
Therefore no \(\operatorname{gamma}(2,\lambda)\) has this
sum law. The source's common-rate requirement cannot be omitted.

If \(X=Z,Y=-Z\) for a standard normal \(Z\), the affine
normal result shows that both marginal laws are standard normal.
Their sum is identically zero and has law \(\delta_0\), with
variance zero. A \(\operatorname{normal}(0,2)\) law has
variance two and no point mass, so it is different. These variables
cannot be independent: their covariance is
\(\mathbb E[-Z^2]=-1\), whereas independent finite-second-moment
variables have zero covariance. This locates the precise failed premise.

Finally, for \(\alpha\leq0\) and \(0<x<1\),
\(x^{\alpha-1}e^{-x}\geq e^{-1}/x\).
The integral of \(1/x\) over \((0,1)\) is infinite: on
each dyadic interval \((2^{-(k+1)},2^{-k}]\), its integral
is at least \(1/2\), and these intervals are disjoint.
Nonnegative countable additivity therefore proves divergence of
the proposed gamma integral at zero. At rate zero the expression
\(\lambda e^{-\lambda x}\) is zero everywhere, with integral
zero rather than one. These calculations justify the stated
excluded parameter cases without extending the normalizing
formulas beyond their proved domains.
\end{proof}
\end{annexitem}
